\documentclass{article}
\usepackage{svg}
\usepackage{amsmath}
\usepackage[utf8]{inputenc}
\usepackage[french]{babel}
\usepackage{graphicx}
\usepackage[T1]{fontenc}
\usepackage{hyperref}
\hypersetup{
    colorlinks,
    citecolor=black,
    filecolor=black,
    linkcolor=black,
    urlcolor=black
}
\graphicspath{images/}

\usepackage{blindtext}

\usepackage{subfiles} % Best loaded last in the preamble

\title{Rapport: Exercice préliminaire}
\author{ }
\date{ }

\begin{document}
\begin{titlepage}
    \begin{center}
        
        {\Large Université de Mons}\\[1ex]
        {\Large Faculté des sciences}\\[1ex]
        {\Large Département d'Informatique}\\[2.5cm]
        
        \newcommand{\HRule}{\rule{\linewidth}{0.3mm}}
        % Title
        \HRule \\[0.3cm]
        { \LARGE \bfseries Stuctures de données II \\[0.3cm]}
        { \LARGE \bfseries Rapport de l'exercice préléminaire \\[0.1cm]} % Commenter si pas besoin
        \HRule \\[1.5cm]
        
        % Author and supervisor
        \begin{minipage}[t]{0.45\textwidth}
            \begin{flushleft} \large
                \emph{Professeur:}\\
                Véronique \textsc{Bruyère}\\
                \emph{Assistant:}\\
                Pierre \textsc{Hauweele}\\
            \end{flushleft}
        \end{minipage}
        \begin{minipage}[t]{0.45\textwidth}
            \begin{flushright} \large
                \emph{Auteurs:} \\
                Theo \textsc{Godin} \\
                Thomas \textsc{Bernard} 
            \end{flushright}
        \end{minipage}\\[2ex]
        
        \vfill
        
        % Bottom of the page
        \begin{center}
            \begin{tabular}[t]{c c c}
                \includegraphics[height=1.5cm]{images/logoumons.jpg} &
                \hspace{0.3cm} &
                \includegraphics[height=1.5cm]{images/logofs.jpg}
            \end{tabular}
        \end{center}~\\
        
        {\large Année académique 2022-2023}
        
    \end{center}
\end{titlepage}

\tableofcontents

\newpage

\section{Question 1}
    \paragraph{Si on ne considère que les coordonnées en x, voyez-vous que celles-ci sont organisées selon une file de priorité ?
    Où est la coordonnée minimum (maximum)? Cette file à priorité est-elle un tas ? Justifiez.}

\section{Question 2}
    \paragraph{Si on ne considère que les coordonnées en y, voyez-vous que celles-ci sont organisées selon un arbre binaire de recherche ?  Où est la coordonnée minimum (maximum)? Justifiez.} 
    Oui on voit bien que les données sont organisées sous forme d'un arbre binaire car une coordonée y inférieur à celle de son père
    deviendra le fils gauche de ce père alors que qu'elle deviendra le fils droit de ce même père dans l'autre cas. La coordonnée minimale se trouve donc dans la feuille la plus à gauche 
    alors que la coordonnée maximale se trouve elle dans la feuille la plus à droite.

\end{document}
